\chapter{Conclusion and Future Works}
\section{Conclusioni}
Questa ricerca ha verificato la capacità di diverse tecniche di Intelligenza Artificiale, quali Linear Regression (LR), Additive Regression (AR), Random SubSpace (RSS), Random Forest (RF), REPTree, modelli M5P, LSTM e GRU, di prevedere il deficit di pressione di vapore (VPD) per la programmazione dell'irrigazione in nove siti in Sicilia relativi alle ex 9 province siciliane (Agrigento, Caltanissetta, Catania, Enna, Messina, Palermo, Ragusa, Siracusa e Trapani). I dati mensili utilizzati per questo studio provengono dal modello di rianalisi atmosferica globale prodotta dall'ECMWF(Centro Europeo per le previsioni metereologiche a medio termine), nello specifico il dataset utilizzato è la rianalisi dei dati medi mensili ERA5 relativi a singoli livelli dal 1958 al 2025.
La scelta di questi input per il modello è attribuita alle loro elevate prestazioni e all'eccellente capacità di apprendimento per relazioni complesse e altamente non lineari, come le serie temporali VPD. I risultati indicano che tutti i modelli sono, più o meno, soddisfacenti nella previsione del VPD. Più specificamente, da un'analisi comparativa, questo studio ha suggerito i modelli AR e RF come quelli più adatti per future ricerche sulla previsione del VPD nelle stazioni studiate.\\
\begin{itemize}
	\item Per quanto riguarda la fase di addestramento in entrambi i tipi di split il Random Forest si attesta come il migliore dei modelli testati con valori che si attestano rispettivamente tra split spaziale e temporale per CC 0,98 e 0,98, MAE 0,05 e 0,05, RMSE 0,08 e 0,07, RAE 15,59\% e 14,93\%, RRSE 18,91\% e 18,10\%, questi valori sono seguiti rispettivamente da quelli di M5P, RSS e AR, come già sottolineato invece il modello REPTree si attesta sempre tra i peggiori anche nella fase di training e le reti neurali ricorrenti hanno valori solidi per CC ma con errori assoluti superiori alla media
	\item Per la fase di test invece il quadro è più equilibrato, come è normale che sia, i valori riscontrati peggiorano per tutti i modelli però ci sono diverse cose utili da sottolineare. Un caso molto interessante è quello del linear regression che per quanto riguarda la fase di test diventa uno dei migliori modelli tra quelli comparati, con risultati comparabili a quelli di Random Forest e Additive Regression che sono gli altri due modelli con risultati più solidi in questa fase, ciò succede perchè a differenza di tutti gli altri modelli, il modello di linear regression mantiene un gap minore tra i risultati della fase di test e quella di train (dove però si avevano dei valori più bassi rispetto alla media), ciò si può facilmente vedere analizzando nello specifico i risultati del modello di Random Forest i quali per esempio nel caso di coefficiente di correlazione peggiorano nella fase di test di quasi 0,4 punti mentre per il linear regression la differenza tra le due fasi è di appena 0,01 punti. Questa differenza tra valorei di test e di train però può dimostrare come i modelli con questo gap possano avere ancora un margine di miglioramento sfuggito anche con la fase di tuning condotta da questo studio.
\end{itemize}
\section{Ipotesi e sviluppi futuri}
\begin{itemize}
	\item Dall'analisi dei risultati si può evincere quindi come i modelli testati riescano tutti a predire con un certo grado di accuratezza i valori di VPD futuri anche a distanza di anni, però si è notato anche come si possa probabilmente ridurre sensibilmente la grandezza del dataset escludendo dallo studio parametri con un basso impatto sulle previsioni di vpd futuri come il valore delle precipitazioni totali o quello dell'evaporazione, che potrebbero nel peggiore dei casi aggiungere rumore e nel migliore semplicemente rallentare il processo di allenamento dei modelli
	\item Un altro elemento da considerare si è riscontrato nell'analisi dei risultati di LSTM e GRU che hanno mostrato dei valori ottimali per quanto riguarda il coefficiente di correlazione e gli errori relativi ma hanno anche prodotto errori assoluti superiori agli altri modelli, ciò potrebbe dimostrare come ci potrebbe essere margine di miglioramento nell'uso di questi modelli quindi questo studio suggerisce come possa dimostrarsi utile il proseguire dell dell'analisi di questa specifica sezione dello studio in futuri lavori, modificando il dataset di partenza utilizzando valori più continui o analizzando più approfonditamente i motivi di questi errori assoluti più alti della media degli altri modelli
	\item infine questo studio si è concentrato su una rianalisi di valori mensili e quindi mediati tra valori rilevati quotidianamente, di conseguenza suggerisce che potrebbe essere utile utilizzare macchine migliori rispetto a quelle utilizzate per l'allenamento dei modelli in questo studio, allo scopo di realizzare uno studio anche sui valori quotidiani di vpd che però, come ipotizzabile, porterebbero alla costruzione di dataset molto più ampi e complessi da gestire. Un'analisi di questo tipo permetterebbe di verificare se l'utilizzo di dati giornalieri possa in qualche modo migliorare le capacità predittive dei modelli analizzati.
\end{itemize}
\label{sec:conclusion}